%%%%%%%%%%  scan idx 116  |  folio 105  %%%%%%%%%%

\medskip
\noindent\textbf{Theorem 4.10.}
Let $X_0$ be an affine scheme of finite type over $k$.

\smallskip
\noindent\textbf{(a)}
If $X_0$ is locally a complete intersection, then
$\underline M_{X_0}$ is smooth over $\underline C_\Lambda$.

\smallskip
\noindent\textbf{(b)}
Suppose that, for a fixed closed point $x\in X_0$, the open
$X_0-\{x\}$ is smooth over $k$. Then both functors
\[
\underline M_{X_0}
\longrightarrow
\underline M_{\operatorname{Spec}(\mathcal O_{X_0,x})}
\longrightarrow
\underline M_{\operatorname{Spec}(\widehat{\mathcal O}_{X_0,x})}
\]
are minimally smooth. In particular, if $(R,\mathfrak X)$ is a formal
versal deformation of $X_0$, then $(R,\mathfrak X_{(x)})$ and
$(R,\mathfrak X_{(\widehat x)})$ are formal versal deformations of
$\operatorname{Spec}(\mathcal O_{X_0,x})$ and
$\operatorname{Spec}(\widehat{\mathcal O}_{X_0,x})$, respectively.

\medskip
\noindent\textit{Proof.}
We use the notation for the cohomology of commutative algebras from
section 3, and write $X_0=\operatorname{Spec}(A_0)$.

%%%%%%%%%%  scan idx 117  |  folio 106  %%%%%%%%%%

\smallskip
\noindent\textbf{(a)}
It suffices to show that
$P:[\underline M_{X_0}]\longrightarrow\underline C_\Lambda$ is
smooth. Let $(R,\operatorname{Spec}(A))$ be a deformation of
$\operatorname{Spec}(A_0)$ over
$R\in\operatorname{Ob}(\underline C_\Lambda)$. Given a small
extension $R'\longrightarrow R$ in $\underline C_\Lambda$, so that
\[
0\longrightarrow k\longrightarrow R'\longrightarrow R
\longrightarrow0
\]
is exact, we must lift $(R,\operatorname{Spec}(A))$ to $R'$. The
transitivity sequence for $D^i$ with coefficients in $A_0$, applied
to $R'\longrightarrow R\longrightarrow A$, gives
\[
\cdots\longrightarrow D^1(A|R,A_0)
\longrightarrow D^1(A|R',A_0)
\longrightarrow D^1(R|R',k)\otimes_kA_0
\longrightarrow D^2(A|R,A_0)\longrightarrow\cdots .
\]
Because $R'\longrightarrow R$ is a small extension,
$D^1(R|R',k)$ is one-dimensional over $k$, generated by the class
$[R']$ represented by
$0\longrightarrow k\longrightarrow R'\longrightarrow R
\longrightarrow0$. Since $A$ is flat over $R$,
\[
D^i(A|R,A_0)=D^i(A_0|k,A_0)
\qquad(i>0).
\]
In particular, $D^2(A|R,A_0)=0$ because $A_0$ is locally a complete
intersection, by the low-degree description in 3.12(8). We therefore
obtain an exact sequence
\[
\longrightarrow D^1(A_0|k,A_0)
\longrightarrow D^1(A|R',A_0)
\longrightarrow D^1(R|R',k)\otimes_kA_0
\longrightarrow0.
\]
Hence there is an $R'$-algebra extension $A'$ of $A$ by $A_0$ whose
class maps to $[R']\otimes1$. Equivalently,
$(R,\operatorname{Spec}(A))$ lifts to the deformation
$(R',\operatorname{Spec}(A'))$.

\smallskip
\noindent\textbf{(b)}
By 3.14 there are isomorphisms
\[
D^1(A_0|k,A_0)
\xrightarrow{\ \sim\ }
D^1((A_0)_x|k,(A_0)_x)
\xrightarrow{\ \sim\ }
D^1(\widehat{(A_0)_x}|k,\widehat{(A_0)_x}),
\]
and hence
\[
[\underline M_{X_0}(k[\varepsilon])]
\xrightarrow{\ \sim\ }
[\underline M_{\operatorname{Spec}(\mathcal O_{X_0,x})}
 (k[\varepsilon])]
\xrightarrow{\ \sim\ }
[\underline M_{\operatorname{Spec}(\widehat{\mathcal O}_{X_0,x})}
 (k[\varepsilon])].
\]
By 2.7(a), it remains to show that
\[
[\underline M_{X_0}]
\longrightarrow
[\underline M_{\operatorname{Spec}(\mathcal O_{X_0,x})}]
\longrightarrow
[\underline M_{\operatorname{Spec}(\widehat{\mathcal O}_{X_0,x})}]
\]
are smooth functors.

Let $(R,\operatorname{Spec}(A))$ be a deformation of
$X_0=\operatorname{Spec}(A_0)$, and let
$0\longrightarrow k\longrightarrow R'\longrightarrow R
\longrightarrow0$ be a small extension. As above, the homomorphisms
$R'\longrightarrow R\longrightarrow A$ give an exact commutative
diagram

%%%%%%%%%%  scan idx 118  |  folio 107  %%%%%%%%%%

\begin{center}
\scriptsize
\begin{tikzcd}[column sep=tiny,row sep=large]
{} \arrow[r]
& D^1(A_0|k,A_0) \arrow[r] \arrow[d,"\sim"']
& D^1(A|R',A_0) \arrow[r,"p"] \arrow[d]
& D^1(R|R',k)\otimes_kA_0 \arrow[r] \arrow[d]
& D^2(A_0|k,A_0) \arrow[r] \arrow[d,"\sim"]
& \cdots \\
{} \arrow[r]
& D^1(A_0|k,(A_0)_x) \arrow[r]
& D^1(A|R',(A_0)_x) \arrow[r]
& D^1(R|R',k)\otimes_k(A_0)_x \arrow[r]
& D^2(A_0|k,(A_0)_x) \arrow[r]
& \cdots \\
{} \arrow[r]
& D^1((A_0)_x|k,(A_0)_x) \arrow[r] \arrow[u]
& D^1(A_x|R',(A_0)_x) \arrow[r,"p_x"] \arrow[u]
& D^1(R|R',k)\otimes_k(A_0)_x \arrow[r] \arrow[u]
& D^2((A_0)_x|k,(A_0)_x) \arrow[r] \arrow[u]
& \cdots
\end{tikzcd}
\end{center}
The indicated isomorphisms follow from the fact that $A_0$ is smooth
over $k$ away from the single closed point $x$, so that
\[
D^i(A_0|k,(A_0)_x)
=D^i(A_0|k,A_0)_x
=D^i(A_0|k,A_0)
\qquad(i>0).
\]

Suppose we are given an arrow
\[
(R',\operatorname{Spec}(B))
\longrightarrow(R,\operatorname{Spec}(A_x))
\]
in $[\underline M_{\operatorname{Spec}((A_0)_x)}]$. This amounts to
a class $[B]\in D^1(A_x|R',(A_0)_x)$ satisfying
$p_x([B])=[R']$. A diagram chase gives a class
$[A']\in D^1(A|R',A_0)$ such that $p([A'])=[R']$ and $A'_x=B$.
Thus $(R',\operatorname{Spec}(A'))$ is a deformation of
$\operatorname{Spec}(A_0)$ over $R'$ whose localization at $x$ is
$\operatorname{Spec}(B)$. This proves that
\[
\underline M_{X_0}\longrightarrow
\underline M_{\operatorname{Spec}(\mathcal O_{X_0,x})}
\]
is smooth.

The smoothness of
\[
\underline M_{\operatorname{Spec}((A_0)_x)}
\longrightarrow
\underline M_{\operatorname{Spec}(\widehat{(A_0)_x})}
\]
follows from the analogous diagram and the identities
\[
D^i((A_0)_x|k,\widehat{(A_0)_x})
=\widehat{D^i((A_0)_x|k,(A_0)_x)}
=D^i((A_0)_x|k,(A_0)_x)
\qquad(i>0),
\]
since $x$ is the only possible nonsmooth point. \hfill$\square$
